\documentclass{article}
\usepackage{amsmath}
\begin{document}
% ===== 编辑器行为演练场：在指定位置按说明操作，不需要编译 =====

\section{Auto-pairing}
% 在下一行行尾依次输入 { [ ( $ ：每个都应自动补闭合并把光标放中间；再输入对应闭合符应“跳过”而不是重复。
Type here:
% 在下一行行尾输入 \{ ：反斜杠后的 { 不应配对。
Escaped:
% 选中下面的单词 WORD 再输入 { ：应变成 {WORD} 且 WORD 保持选中。
Wrap WORD with braces.
% 在 {} 中间按 Delete：应一次删掉一对。
Delete pair: {}

\section{Environments}
% 在下一行行尾输入 \begin{itemize} 然后回车：应自动补 \end{itemize}，光标在中间缩进两格。
Type begin here:
% 下面已经有 \end{enumerate}：在 \begin 行末回车不应重复补一个 \end。
\begin{enumerate}
\end{enumerate}
% 输入 \begin{ 后应自动弹出环境候选；选 align 后同样补 \end{align}。
Completion:

\section{Indent and comment}
% 选中下面三行按 Tab / Shift-Tab（或 Cmd-] Cmd-[），再按 Cmd-/ 两次。
line one
  line two (already indented)
line three
% 单行无选区按 Tab 插入两个空格；Shift-Tab 反缩进。
    four-space line

\section{Enter keeps indentation}
    % 在这一行末尾回车：新行应保留 4 个空格缩进。

\section{Completion}
% 在下面每个 { 里按 Ctrl-Space。
\ref{}  \cite{}  \eqref{}  \pageref{}  \autoref{}
% 输入 \sec 后 Ctrl-Space：应列出 \section \subsection 等；\vect 是本文件定义的命令，也应在列。
\newcommand{\vect}[1]{\mathbf{#1}}
Command:
\label{sec:first}\label{fig:second}\label{tab:third}

\section{Bracket matching}
% 把光标放在每个括号旁边，看配对括号是否高亮；转义的 \{ 不应匹配。
\frac{a}{b} [x] (y) \{ not paired \} \left( \frac{1}{2} \right) {outer {inner} outer}
% 不平衡的括号：{ 没有闭合 ------ 光标放这里不应卡顿（20000 字符扫描上限）。

\section{Go to line and sync}
% Cmd-L 输入 40；再输入 9999（超出范围应跳到末尾或不动，不崩溃）。
% 在任意行 Cmd-J 而 PDF 尚未编译：状态栏应提示“还没有 PDF”。

\section{Find and replace}
% Cmd-F 查找 needle：3 处；Opt-Cmd-F 全部替换成 pin；Cmd-Z 撤销。
needle needle needle
% Cmd-E 用所选内容查找：选中 haystack 再 Cmd-E Cmd-G。
haystack

\end{document}
